\section{Introduction}
\label{sec:introduction}

\IEEEPARstart{M}{ulti-robot} systems are assembled from components whose
correctness claims rest on different assumptions.  A mission synthesizer may
guarantee a temporal goal within a declared context; a planner may report
feasibility for a particular state and model; a learned predictor may forecast
peer motion over a horizon; middleware may report a quality-of-service event;
and a safety filter may guarantee forward invariance under a dynamics and timing
model.  Mission success depends on each claim and on whether its assumptions
still hold when another component, or another robot, consumes its output.  The
practical requirement is easy to state: \emph{a result should cross an
interface with enough context for the receiver to decide whether it is still
valid.}

This requirement is becoming harder to meet.  Multi-robot stacks now combine
classical task allocation and motion coordination
\cite{gerkey2004mrta,choi2009cbba,sharon2015cbs} with language-conditioned
planning and learned prediction \cite{mandi2024roco,obata2025lipllm,li2026llmmrs}.
Learned components expand capability, but they do not remove the need for
explicit physical, temporal, and deployment assumptions.  Inference is also
distributed: a prediction may be computed on board, next to the controller that
uses it, or remotely on infrastructure that fuses more data at the cost of
latency \cite{kehoe2015cloud}.  Each placement changes the interfaces that a
result must cross.  At lower layers, ROS~2 and the Data Distribution Service
(DDS) transport state and commands, while runtime-assurance and barrier-based
controllers may override a nominal decision before actuation
\cite{macenski2022ros2,omg2015dds,desai2019soter,wang2017barrier}.  An
end-to-end execution argument must therefore cross semantic, physical,
software, communication, and safety boundaries.

Multi-robot coupling makes this more than an ordinary integration problem.
Stale or ambiguous state at robot $r_i$ can make a locally admissible decision
invalidate robot $r_j$'s route, resource claim, or safety margin.  Join, leave,
restart, and partition events can invalidate assignments even when every message
is syntactically well formed.  The relevant interfaces are \emph{vertical},
between functions within one robot, and \emph{horizontal}, between robots or
distributed services (\cref{fig:command-path}).

\begin{figure}[htbp]
  \centering
  \begin{tikzpicture}[
    x=1cm,y=1cm,
    stage/.style={draw,rounded corners=2pt,minimum height=8mm,
      minimum width=46mm,align=center,font=\small,inner sep=3pt},
    lane/.style={stage,minimum width=42mm},
    flow/.style={-{Latex[length=2.2mm]},line width=0.8pt},
    peer/.style={{Latex[length=2.4mm]}-{Latex[length=2.4mm]},line width=1.1pt,
      draw=MidnightBlue},
    tag/.style={midway,fill=white,inner xsep=3pt,inner ysep=1.5pt,
      font=\scriptsize,align=center}
  ]
    \definecolor{IntentGray}{HTML}{E6E8EB}
    \definecolor{SemanticBlue}{HTML}{DCEBFA}
    \definecolor{PredictPurple}{HTML}{E9E0F5}
    \definecolor{PhysicsGreen}{HTML}{DDEFD8}
    \definecolor{RuntimeOrange}{HTML}{F9E4C8}
    \definecolor{SafetyRed}{HTML}{F7D8D5}

    \node[stage,fill=IntentGray] (intent) at (0,0) {Operator intent};
    \node[stage,fill=SemanticBlue] (task) at (0,-1.25)
      {Semantic task formation};
    \node[stage,fill=PredictPurple] (predict) at (0,-2.50)
      {Prediction of team and world state};
    \node[stage,fill=IntentGray] (joint) at (0,-3.75)
      {Joint allocation, constraints, and forecast};

    \draw[flow] (intent) -- node[tag] {$C_0$: interpreted request} (task);
    \draw[flow] (task) -- node[tag] {$C_1$: structured task} (predict);
    \draw[flow] (predict) -- node[tag] {$C_2$: predicted team state} (joint);

    \foreach \side/\xpos/\idx in {L/-4.6/1,R/4.6/2} {
      \node[lane,fill=PhysicsGreen] (ground\side) at (\xpos,-5.55)
        {Physical grounding};
      \node[lane,fill=RuntimeOrange] (runtime\side) at (\xpos,-6.85)
        {Distributed runtime};
      \node[lane,fill=SafetyRed] (safety\side) at (\xpos,-8.15)
        {Runtime safety and authority};
      \node[lane,fill=IntentGray] (robot\side) at (\xpos,-9.45)
        {Physical robot $r_{\idx}$};
      \draw[flow] (ground\side) -- node[tag] {vertical interface
        $e^{\mathrm v}$} (runtime\side);
      \draw[flow] (runtime\side) -- (safety\side);
      \draw[flow] (safety\side) -- (robot\side);
    }

    \coordinate (split) at (0,-4.45);
    \draw[line width=0.8pt] (joint.south) -- (split);
    \draw[flow] (split) -| (groundL.north);
    \draw[flow] (split) -| (groundR.north);
    \node[font=\small\bfseries,fill=white,inner xsep=3pt] at (-4.6,-4.85)
      {Robot $r_1$};
    \node[font=\small\bfseries,fill=white,inner xsep=3pt] at (4.6,-4.85)
      {Robot $r_2$};

    \draw[peer] (runtimeL.east) -- node[tag,text=MidnightBlue,
      font=\scriptsize\bfseries] {horizontal interface $e^{\mathrm h}$\\
      peer state and coordination} (runtimeR.west);

    \node[draw,dashed,rounded corners=2pt,align=center,font=\footnotesize,
      inner sep=4pt,text width=150mm] at (0,-10.55)
      {Each admitted transfer carries a typed payload together with semantic
       and physical, temporal, uncertainty and evidence, compatibility and
       lifecycle, and authority and recovery obligations.};
  \end{tikzpicture}
  \caption{Cross-layer command path of a two-robot team.  Single-headed arrows
  are producer--consumer transfers; $C_0$--$C_2$ label the shared
  command-formulation interfaces.  Vertical interfaces connect functions within
  one robot, and the double-headed horizontal interface exchanges peer state and
  coordination obligations between robots.}
  \label{fig:command-path}
\end{figure}


Existing work addresses parts of this problem through contract-based mission
design, quantitative mission assurance, component contracts for the Robot
Operating System (ROS), compositional reasoning for interconnected dynamics,
runtime monitoring, and safety-controller switching
\cite{mallozzi2020crome,chaki2013assurance,luckcuck2026compositional,
sharf2021dynamical,ghaffari2024rosmonitoring,desai2019soter}.  These approaches
protect different objects and produce different kinds of evidence.  This paper
asks whether their combined coverage supports the full multi-robot
command-to-action path:
\begin{enumerate}[\IEEEsetlabelwidth{RQ3:}]
  \item[RQ1:] Which objects, assumptions, and guarantees do existing
    contract-based and runtime-assurance approaches represent?
  \item[RQ2:] How are these claims composed, monitored, enforced, or repaired,
    and what evidence supports them?
  \item[RQ3:] Which obligations remain implicit when information crosses
    semantic, physical, distributed, and inter-robot boundaries, and what must
    an interface carry to make them checkable?
\end{enumerate}

The paper makes four contributions:
\begin{enumerate}
  \item a cross-layer taxonomy of contract-based and runtime-assurance
    approaches, organized by protected object, mechanism, and reported evidence
    (\cref{sec:taxonomy});
  \item a coverage analysis that identifies five recurring interface gaps,
    supported by a qualitative coverage matrix and by public failure reports
    from ROS~2 middleware (\cref{sec:gaps});
  \item an interface-contract reference framework with three-valued admission
    and recovery semantics, including a temporal admission margin for peer
    state that determines which metadata a horizontal interface must carry and
    when on-board inference is preferable to remote inference
    (\cref{sec:framework}); and
  \item a failure-attribution taxonomy and a worked command trace that assign
    validators and recovery authority along the command path
    (\cref{sec:attribution,sec:trace}).
\end{enumerate}

The framework is a reference model.  It does not define a new contract algebra,
and it has not been evaluated as a reliability mechanism.
\Cref{sec:background} states the interface problem, \cref{sec:method} describes
the review method, \cref{sec:taxonomy} presents the taxonomy, and
\cref{sec:gaps} derives the gaps.  \Cref{sec:framework} develops the
framework, \cref{sec:discussion} discusses its implications and limitations,
and \cref{sec:conclusion} concludes.
